% Generated draft: requires mathematical and source review.
\documentclass[11pt]{article}
\usepackage[margin=25mm]{geometry}
\usepackage{fontspec}
\setmainfont{Latin Modern Roman}
\usepackage{amsmath,amssymb,mathtools,unicode-math}
\setmathfont{Latin Modern Math}
\usepackage{enumitem,needspace,xurl,hyperref}
\setlength{\parindent}{0pt}
\setlength{\parskip}{5pt}
\setlength{\emergencystretch}{3em}
\newcommand{\N}{\mathbb{N}}
\newcommand{\Z}{\mathbb{Z}}
\newcommand{\Q}{\mathbb{Q}}
\newcommand{\R}{\mathbb{R}}
\newcommand{\C}{\mathbb{C}}
\newcommand{\F}{\mathbb{F}}
\newcommand{\A}{\mathbb{A}}
\newcommand{\PP}{\mathbb{P}}
\newcommand{\E}{\mathbb{E}}
\newcommand{\eps}{\varepsilon}
\newcommand{\dd}{\,\mathrm{d}}
\begin{document}
\section*{UnsolvedMath ID: 30006557}
\textbf{Source problem code:} catalogue entry
\subsection{Definitions and mathematical statement}
For the usual finite simple undirected graph interpretation, a plane drawing of a graph \(G\) assigns distinct points to its vertices and simple arcs to its edges, with the prescribed endpoints and with no other vertex on an edge. Work in general position: intersections of edge interiors are finitely many transverse crossings, with no three edge interiors meeting at one point. The number of crossings counts intersection points, not merely pairs of intersecting edges. Define the crossing number \(\operatorname{cr}(G)\) as the minimum number of crossings over such drawings.

A drawing is \(x\)-monotone if every vertical line intersects each edge arc in at most one point. Define the monotone crossing number \(\operatorname{mon\text{-}cr}(G)\) as the minimum number of crossings over \(x\)-monotone drawings. Vertex positions and their horizontal order are allowed to vary in this minimum; the supplied record prescribes no vertex order.

The supplied question asks whether there exists an absolute real constant \(C>0\), independent of \(G\), such that
\[
\operatorname{mon\text{-}cr}(G)\leq C\,\operatorname{cr}(G)
\qquad\text{for every graph }G.
\]
The quantifier includes graphs with \(\operatorname{cr}(G)=0\); no additive error term or positive-crossing hypothesis is supplied.

\emph{Unresolved scope qualification.} The record does not explicitly specify whether graphs are finite, simple, or undirected, nor does it supply the drawing conventions defining monotone crossing number. The definitions above state the usual finite simple undirected interpretation explicitly. Their agreement with the report's intended graph class and conventions requires source review; no assertion about loops, parallel edges, infinite graphs, or prescribed vertex orders is inferred.
\subsection{Short English statement}
Can the minimum number of crossings in a drawing with horizontally monotone edges always be bounded by one universal constant times the ordinary crossing number? The supplied record labels this question open. The intended graph class and drawing conventions require confirmation.
\subsection{Sources}
\begin{itemize}
\item Local UnsolvedMath registry, permanent ID \(30006557\), supplied complete source record, title: Comparing Monotone and Ordinary Crossing Numbers. Exact source problem code: \texttt{catalogue entry}. The record gives proposed year \(2026\) and names no proposer.
\item Supplied report citation: \emph{Combinatorics} (2026), Oberwolfach Reports, pp.\ 5--79, DOI \texttt{10.4171/owr/2026/1}. \url{https://doi.org/10.4171/owr/2026/1}. The supplied background describes Oberwolfach Reports as documenting mathematical workshops and their open-problem discussions. This reference is also listed among the supplied literature sources as Oberwolfach Report 1/2026.
\item Supplied literature reference: Pach--Toth, \emph{Monotone Crossing Number}. \url{https://doi.org/10.1007/978-3-642-25878-7_27}. No further bibliographic details are supplied.
\item Literature assessment dated exactly \texttt{2026-08-22}: the supplied curation labels the problem open and states that the best general comparison gives a quadratic-type upper bound for monotone crossing number in terms of crossing number, while no universal linear constant is known. This is supplied dated curation derived from the OWR material, not an independently verified literature assessment. Its time-sensitive status and cited claims require independent verification; no explicit quadratic inequality or hypotheses are supplied.
\item The registry assigns statement status \texttt{corrected\_verified} and reports that the official Oberwolfach report was checked, the question isolated from surrounding exposition, and its notation repaired. These are supplied provenance claims, not verification performed for this draft. The supplied original-statement field also says that a report passage is retained verbatim as context, but no separately identified surrounding passage is provided.
\end{itemize}
\end{document}
